\section{Fourier Transform - Theory}

\subsection{Fourier Series vs Transform}

The Fourier series represents a \textit{periodic} signal as a discrete linear combination of 
sine and cosine functions at integer multiples of a fundamental frequency.
The Fourier transform generalizes this to \textit{non-periodic} signals by projecting 
the signal onto complex exponentials across a continuum of frequencies,
replacing the discrete sum with an integral.

\subsection{Derivation}

We will now show the derivation of the Fourier transform of $f(x):=\delta(x-d)+\delta(x+d)$
\begin{proof}
If we examine the Fourier transform simply by inserting $f(x)$ in the definition, we see:
\begin{equation}
	\mathcal{F}\{f\}(u) 
	= F(u)
	= \int_{-\infty}^{\infty} \big(
		\delta(x-d)+\delta(x+d)
	\big)e^{-i 2 \pi u x} dx
\end{equation}

And by definition of the dirac-delta function, we know that the expression
will be 0 everywhere except at $x=d$ and $x=-d$.
By first rewriting by linearity of the integral, we can then evaluate the integral at the two points:
\begin{equation}
	F(u)
	= \int_{-\infty}^{\infty} \delta(x-d)e^{-i 2 \pi u x} dx
	+ \int_{-\infty}^{\infty} \delta(x+d)e^{-i 2 \pi u x} dx
\end{equation}

Applying a shift translation by $d$ and $-d$ for the first and second integral respectively,
we get them on the following form:

\begin{equation}
	F(u)
	= \int_{-\infty}^{\infty} \delta(x)e^{-i 2 \pi u (x+d)} dx
	+ \int_{-\infty}^{\infty} \delta(x)e^{-i 2 \pi u (x-d)} dx
\end{equation}
And by the shifting property of the dirac-delta function,
which states
$\int_{-\infty}^{\infty} \delta(x-a)f(x) dx = f(a)$,
\footnote{
	From lecture slides on Convolution week 1
}
we can evaluate the two integrals as follows:
\begin{equation}
	F(u)
	= \int_{-\infty}^{\infty} \delta(x)e^{-i 2 \pi u (x+d)} dx
	+ \int_{-\infty}^{\infty} \delta(x)e^{-i 2 \pi u (x-d)} dx
	= e^{i 2 \pi u d} + e^{-i 2 \pi  u d}.
\end{equation}
Finally by Euler's formula, we can rewrite this as:
\begin{align}
	F(u) 
	=& \big(
		\cos(2\pi u d) + i \sin(2\pi u d)
	\big) + \big(
		\cos(-2\pi u d) + i \sin(-2\pi u d)
	\big) \\
	=&  \big(
		\cos(2\pi u d) + i \sin(2\pi u d)
	\big) + \big(
		\cos(2\pi u d) - i \sin(2\pi u d)
	\big) \\
	=& 2 \cos(2\pi u d)
\end{align}
Where the intermediate step (2) follows from the fact that $\cos$ is an even function and $\sin$ is an odd function.
\end{proof}
